\documentclass[doctype=article,language=arabic]{../../omnilatex}

\title{دعم اللغة العربية في أومني لاتك}
\subtitle{مثال توضيحي للتخطيط من اليمين إلى اليسار}
\author{أحمد محمد}
\date{\today}

\begin{document}
\maketitle

\begin{abstract}
يُقدم هذا المستند مثالاً توضيحياً على دعم اللغة العربية في نظام
أومني لاتك للتنضيد. يتضمن تخطيطاً من اليمين إلى اليسار مع الخطوط
العربية والإعدادات المناسبة للغة.
\end{abstract}

\section{مقدمة}

اللغة العربية هي إحدى أكثر اللغات انتشاراً في العالم، إذ يتحدث بها
أكثر من أربعمائة مليون شخص. تتميز الكتابة العربية بالاتجاه من اليمين
إلى اليسار، مما يتطلب معالجة خاصة في أنظمة التنضيد الإلكتروني.
يُوفر نظام أومني لاتك دعماً كاملاً لهذا الاتجاه من خلال حزمة
التنضيد ثنائي الاتجاه.

\section{الرياضيات في النص العربي}

يمكن كتابة المعادلات الرياضية داخل النص العربي مع الحفاظ على
الاتجاه الصحيح من اليسار إلى اليمين. على سبيل المثال، معادلة
الطاقة الكهربائية:

\begin{equation}
  E = mc^{2}
  \label{eq:energy}
\end{equation}

وكذلك المعادلة التربيعية:

\begin{equation}
  x = \frac{-b \pm \sqrt{b^{2} - 4ac}}{2a}
  \label{eq:quadratic}
\end{equation}

\section{القوائم}

\begin{itemize}
  \item العنصر الأول في القائمة
  \item العنصر الثاني
  \item العنصر الثالث مع شرح إضافي
\end{itemize}

\begin{enumerate}
  \item الخطوة الأولى
  \item الخطوة الثانية
  \item الخطوة الثالثة
\end{enumerate}

\section{الجداول}

\begin{table}[htbp]
  \caption{مقارنة بين أنظمة التنضيد}
  \begin{tabular}{llr}
    \textbf{النظام} & \textbf{النوع} & \textbf{السنة} \\
    \hline
    تِي‌إكس & مفتوح المصدر & ١٩٧٨ \\
    لاتك & مفتوح المصدر & ١٩٨٤ \\
    أومني لاتك & مفتوح المصدر & ٢٠٢٤ \\
  \end{tabular}
\end{table}

\section{الخاتمة}

يُثبت هذا المثال أن نظام أومني لاتك قادر على التعامل مع النصوص
العربية بكفاءة عالية، مع الحفاظ على جودة التنضيد ودعم المعادلات
الرياضية والجداول والقوائم بالاتجاه الصحيح.

\end{document}
